\chapter{Glossary}
\label{annexe:glossaire}

\begin{description}
  \item[Proxmox VE] Open-source virtualisation platform (virtual machines and LXC containers).
  \item[Docker / Docker Compose] Containerisation and multi-container orchestration tools running on a single machine.
  \item[Traefik] Reverse proxy that routes HTTP(S) traffic to containers and manages TLS certificates automatically.
  \item[Let's Encrypt] Free certificate authority issuing TLS certificates.
  \item[CGNAT (Carrier-Grade NAT)] Network address translation applied by an ISP across many customers sharing the same public IPv4 address, which makes inbound port forwarding to a specific customer's equipment impossible.
  \item[Cloudflare Tunnel] An agent (\texttt{cloudflared}) that opens an outbound connection from a server to Cloudflare, allowing public HTTPS traffic to reach it without any inbound port being opened -- works behind CGNAT.
  \item[DuckDNS] Free dynamic DNS service.
  \item[Jenkins] Continuous integration and continuous deployment (CI/CD) server.
  \item[UFW] Simplified firewall front-end for Linux (built on iptables).
  \item[Fail2ban] Service that temporarily bans an IP address after repeated suspicious login attempts.
  \item[cloud-init] Cloud instance initialisation tool; on Ubuntu it can drop configuration files (e.g. under \texttt{/etc/ssh/sshd\_config.d/}) that silently override manual edits to the main configuration file.
  \item[Webhook] Automatic HTTP notification sent by a service to a URL whenever a given event occurs.
  % TODO: extend this list as new technical terms appear while writing.
\end{description}
